% !PRINTED 627
\weberpage{627}
b) If $p-1$ is not divisible by 4, and hence
\weq{11}{p \equiv -5^\beta \, (\mathrm{mod}\ m),}
then $f$ is the least positive number that satisfies the congruences
\[ f \equiv 0 \, (\mathrm{mod}\ 2), \quad 2\beta f \equiv 0 \, [\mathrm{mod}\ \varphi(m)] \]
and consequently $e$ is the greatest common divisor of $\varphi(m)$ and $2\beta$, or, if $2\beta$ is divisible by $\varphi(m)$, equal to $\tfrac{1}{2}\varphi(m)$. (An exception is formed here by the case $m=4$, in which $e=1$. This case, which has indeed already been completely settled in earlier chapters, we therefore now set aside.)

It then follows from (11) that a number of the form $5^\lambda$ can be congruent to a power of $p$ only if $\lambda$ is a multiple of $2\beta$, and hence a multiple of $e$.

Now here a number can be found in each of the subsidiary groups $\mathfrak{A}\xi$ that is $\equiv 1 \, (\mathrm{mod}\ 4)$, since $p \equiv -1 \, (\mathrm{mod}\ 4)$ is a number in $\mathfrak{A}$, and we accordingly obtain a complete system of the $\xi$ in the form:
\[ 1, 5, 5^2, \ldots, 5^{e-1}, \]
and in it $e$ is at most $= \tfrac{1}{2}\varphi(m) = 2^{\varkappa-2}$.

We collect what has been proved in the following theorem:

III. \gesperrt{If one decomposes the prime $p$ different from $q$ in the field $\Omega_m$ into its prime factors:}
\[ p = \mathfrak{p}_{\xi_1} \mathfrak{p}_{\xi_2} \cdots \mathfrak{p}_{\xi_e}, \]
\gesperrt{then, if $p-1$ is divisible by $q$, or in the case of an even $m$ by 4, one may let the $\xi$ run through a complete system of residues not divisible by $q$ with respect to the modulus $m_1$, and take the $\mathfrak{p}_\xi$ to be prime functionals in the field $\Omega_{m_1}$, if $m_1$ is the greatest common divisor of $p-1$ and $m$.}

\gesperrt{But if $p-1$ is not divisible by $q$, or for an even $m$ not by 4, and if $g$ denotes in the first case a primitive root of $m$, in the second case the number 5, then $\xi$ runs through the sequence of numbers}
\[ 1, g, g^2, \ldots, g^{e-1}, \]
\gesperrt{and for an odd $m$, $e$ is not divisible by $q-1$, and for an even $m$ it is at least $=2$ and at most $=\tfrac{1}{2}\varphi(m)$. Excluded from this is the case $m=4$.}
